\chapter{HTTP API Reference}

This appendix is a desk reference. You do not need to read it from top to bottom before using the project. Return here when you are writing a request, checking a field name, or wondering which service owns an address. The story behind these calls lives in Chapters 3 through 9.

\section{Conventions}

The dashboard uses JSON for ordinary requests and responses, multipart form data for project upload, server-sent events for CSV progress, HTML for profiling reports, and raw octet streams for firmware downloads.

OTA failures share this envelope:

\begin{lstlisting}[language=JSON,numbers=none]
{
  "success": false,
  "error": {
    "code": "TARGET_MISMATCH",
    "message": "Device chip does not match firmware target",
    "details": "optional diagnostic detail",
    "hint": "optional corrective action"
  }
}
\end{lstlisting}

Fields with no value may be omitted. IDs are UUID strings except \codepath{device_id}, which is chosen by configuration or derived from the station MAC.

\section{OTA endpoint index}

{\small
\begin{tabularx}{\textwidth}{>{\raggedright\arraybackslash}p{0.1\textwidth}>{\raggedright\arraybackslash}p{0.52\textwidth}Y}
\toprule
Method & Endpoint & Purpose\\
\midrule
GET & \codepath{/api/ota/status} & Service health and public URL\\
GET & \codepath{/api/ota/projects} & List uploaded projects\\
POST & \codepath{/api/ota/projects} & Upload a project ZIP\\
POST & \codepath{/api/ota/projects/<project-id>/build} & Queue a build\\
GET & \codepath{/api/ota/builds} & List build history\\
GET & \codepath{/api/ota/builds/<build-id>} & Build status and full log\\
GET & \codepath{/api/ota/firmware} & List ready firmware\\
GET & \codepath{/api/ota/firmware/latest?target=esp32s3} & Newest ready target image\\
GET & \codepath{/api/ota/firmware/<firmware-id>} & Firmware details\\
GET & \codepath{/api/ota/firmware/<firmware-id>/download} & Stream application image\\
DELETE & \codepath{/api/ota/firmware/<firmware-id>} & Delete unreferenced image\\
GET & \codepath{/api/ota/devices} & List registered devices\\
POST & \codepath{/api/ota/devices/register} & Create or update a device\\
POST & \codepath{/api/ota/devices/<device-id>/heartbeat} & Refresh device state\\
POST & \codepath{/api/ota/devices/<device-id>/deploy} & Select desired firmware\\
GET & \codepath{/api/ota/devices/<device-id>/update} & Poll desired firmware\\
POST & \codepath{/api/ota/devices/<device-id>/ota-status} & Report OTA progress/result\\
\bottomrule
\end{tabularx}
}

\section{Status}

\textbf{GET \codepath{/api/ota/status}} returns 200.

\begin{lstlisting}[language=JSON,numbers=none]
{
  "service": "CSV_PROFILING OTA Server",
  "language": "rust",
  "status": "online",
  "version": "0.1.0",
  "public_base_url": "http://192.168.1.5:7000",
  "supported_targets": ["esp32s3"]
}
\end{lstlisting}

\section{Projects and builds}

\textbf{POST \codepath{/api/ota/projects}} expects multipart fields \codepath{project}, \codepath{target}, and \codepath{version}. A successful upload returns 201.

\begin{lstlisting}[numbers=none]
curl -F "project=@uploaded-firmware-basic.zip" \
  -F "target=esp32s3" -F "version=1.0.1" \
  http://localhost:7000/api/ota/projects
\end{lstlisting}

\begin{lstlisting}[language=JSON,numbers=none]
{
  "project_id": "55e4...",
  "name": "uploaded-firmware-basic",
  "target": "esp32s3",
  "version": "1.0.1",
  "status": "uploaded",
  "project_type": "managed_application",
  "application_api_version": 2,
  "entrypoint": "src/main.rs"
}
\end{lstlisting}

Standalone uploads omit the API version and entrypoint and report \codepath{project_type} as \codepath{standalone}.

\textbf{GET \codepath{/api/ota/projects}} returns an array of project records. Internal filesystem paths are skipped during serialization.

\textbf{POST \codepath{/api/ota/projects/<project-id>/build}} returns 202:

\begin{lstlisting}[language=JSON,numbers=none]
{"build_id":"a8c4...","status":"building"}
\end{lstlisting}

The durable row begins queued and is moved to building by the coordinator task. Only one queued or active build is allowed for a project.

\textbf{GET \codepath{/api/ota/builds}} returns all builds, newest first. \textbf{GET \codepath{/api/ota/builds/<build-id>}} returns one response shaped like:

\begin{lstlisting}[language=JSON,numbers=none]
{
  "build_id": "a8c4...",
  "project_id": "55e4...",
  "project_name": "uploaded-firmware-basic",
  "target": "esp32s3",
  "version": "1.0.1",
  "status": "success",
  "logs": ["Build queued", "Build started", "Tooling: cargo ..."],
  "exit_code": 0,
  "error": null,
  "started_at": "2026-09-07T10:00:00Z",
  "finished_at": "2026-09-07T10:02:00Z",
  "created_at": "2026-09-07T09:59:59Z"
}
\end{lstlisting}

Failed responses may include \codepath{hint}.

\section{Firmware}

List, detail, and latest routes return the same firmware representation:

\begin{lstlisting}[language=JSON,numbers=none]
{
  "firmware_id": "c53b...",
  "project_id": "55e4...",
  "build_id": "a8c4...",
  "project": "uploaded-firmware-basic",
  "target": "esp32s3",
  "version": "1.0.1",
  "filename": "firmware-uploaded-firmware-basic-v1.0.1-a8c4.bin",
  "size": 945332,
  "sha256": "f2534c...",
  "status": "ready",
  "created_at": "2026-09-07T10:02:00Z",
  "download_url": "http://192.168.1.5:7000/api/ota/firmware/c53b.../download"
}
\end{lstlisting}

\textbf{GET \codepath{/download}} returns \codepath{application/octet-stream}, content length, attachment disposition, and \codepath{X-Firmware-SHA256}. The route reconstructs the path from the configured build directory and stored filename, then canonicalizes it before opening.

\textbf{DELETE \codepath{/api/ota/firmware/<firmware-id>}} returns 204. It returns 409 when a device or deployment references the image.

\section{Devices and deployments}

\textbf{POST \codepath{/api/ota/devices/register}} returns 201 and the current device record.

\begin{lstlisting}[language=JSON,numbers=none]
{
  "device_id": "ESP32-S3-001",
  "name": "Sensor Node 1",
  "chip": "esp32s3",
  "current_version": "1.0.0",
  "ip_address": "192.168.1.24"
}
\end{lstlisting}

\textbf{POST \codepath{/api/ota/devices/ESP32-S3-001/heartbeat}} returns 200 and the updated device.

\begin{lstlisting}[language=JSON,numbers=none]
{
  "current_version": "1.0.0",
  "status": "online",
  "ip_address": "192.168.1.24"
}
\end{lstlisting}

Allowed heartbeat statuses are online, offline, updating, rebooting, updated, failed, and unknown.

\textbf{GET \codepath{/api/ota/devices}} returns records with these public fields:

\begin{lstlisting}[language=JSON,numbers=none]
{
  "device_id": "ESP32-S3-001",
  "name": "Sensor Node 1",
  "chip": "esp32s3",
  "ip_address": "192.168.1.24",
  "current_version": "1.0.0",
  "desired_version": "1.0.1",
  "desired_firmware_id": "c53b...",
  "status": "online",
  "ota_status": "pending",
  "ota_progress": 0,
  "ota_error": null,
  "last_seen": "2026-09-07T10:03:00Z"
}
\end{lstlisting}

\textbf{POST \codepath{/api/ota/devices/<device-id>/deploy}} expects a firmware ID and returns 202.

\begin{lstlisting}[language=JSON,numbers=none]
{"firmware_id":"c53b..."}
\end{lstlisting}

\begin{lstlisting}[language=JSON,numbers=none]
{
  "deployment_id": "d911...",
  "device_id": "ESP32-S3-001",
  "firmware_id": "c53b...",
  "desired_version": "1.0.1",
  "status": "pending"
}
\end{lstlisting}

\textbf{GET \codepath{/api/ota/devices/<device-id>/update}} returns either \codepath{update_available: false} or the manifest shown in Chapter 9.

\textbf{POST \codepath{/api/ota/devices/<device-id>/ota-status}} accepts:

\begin{lstlisting}[language=JSON,numbers=none]
{
  "status": "downloading",
  "progress": 35,
  "current_version": null,
  "error": null
}
\end{lstlisting}

Allowed OTA states are idle, pending, downloading, verifying, installing, rebooting, success, and failed. Progress must be between 0 and 100. Failed requires an error. Success defaults to 100 and clears desired firmware after updating the current version.

\section{CSV endpoint index}

These routes belong to Flask on port 5000.

\begin{longtable}{p{0.13\textwidth}p{0.38\textwidth}p{0.39\textwidth}}
\toprule
Method & Endpoint & Purpose\\
\midrule
\endhead
GET & \codepath{/} & Serve the dashboard HTML\\
POST & \codepath{/upload} & Accept multipart CSV field \codepath{file} and return \codepath{job_id}\\
GET & \codepath{/progress/<job-id>} & Stream profiling progress with SSE\\
GET & \codepath{/report/<job-id>} & Serve generated report HTML\\
GET & \codepath{/sessions} & List retained report sessions\\
GET & \codepath{/session/<job-id>} & Check and restore one session\\
POST & \codepath{/clean/<job-id>} & Apply cleaning controls and refresh metadata\\
POST & \codepath{/playground/<job-id>} & Run preview, analysis, chart, or explicit mutation\\
GET & \codepath{/download-cleaned/<job-id>} & Download cleaned CSV\\
DELETE & \codepath{/reset/<job-id>} & Delete one CSV session\\
DELETE & \codepath{/reset-all} & Delete all CSV sessions\\
\bottomrule
\end{longtable}

Playground request fields include \codepath{operation}, \codepath{mutate}, \codepath{limit}, \codepath{column}, \codepath{column_2}, \codepath{value}, \codepath{operator}, \codepath{indexes}, and \codepath{chart_type}. The endpoint clamps limit to 1 through 500.
